
\documentclass[12pt,a4paper]{article}
\usepackage[utf8]{inputenc}
\usepackage[T1]{fontenc}
\usepackage[spanish]{babel}
\usepackage{amsmath,amssymb,amsfonts}
\usepackage{graphicx}
\usepackage{geometry}
\usepackage{mathpazo}        % Fuente Palatino
\usepackage{bm}              % Negrita para símbolos matemáticos
\linespread{1.05}            % Espaciado ampliado
\setlength{\parskip}{0.5em}  % Espacio entre párrafos
\geometry{margin=2.5cm}

\title{1\quad Ecuaciones Básicas}
\author{Traducción técnica del texto de Feistauer (2006)}
\date{}

\begin{document}

\maketitle

\section*{1\quad Ecuaciones Básicas}

Los problemas de dinámica de fluidos juegan un papel importante en muchas áreas de la ciencia y la tecnología. Por ejemplo:
\begin{itemize}
  \item la industria aeronáutica,
  \item la ingeniería mecánica,
  \item la turbomaquinaria,
  \item la industria naval,
  \item la ingeniería civil,
  \item la ingeniería química,
  \item la industria alimentaria,
  \item la protección ambiental,
  \item la meteorología,
  \item la oceanografía,
  \item la medicina.
\end{itemize}

La visualización del flujo puede obtenerse mediante:
\begin{enumerate}
  \item experimentos (por ejemplo, en túneles de viento), que son costosos, prolongados o incluso inviables;
  \item modelos matemáticos, resueltos numéricamente con computadoras modernas. Este enfoque constituye la \textit{Dinámica de Fluidos Computacional} (CFD), cuyo objetivo es obtener resultados comparables a los de los experimentos físicos y, cuando sea posible, reemplazarlos.
\end{enumerate}

\section*{1.1\quad Ecuaciones Gobernantes}

Sea $\Omega \subset \mathbb{R}^N$ con $N=1,2,3$ el dominio ocupado por el fluido, y sea $\mbox{$(0,T) \subset \mathbb{R}$}$ un intervalo temporal en el cual observamos el movimiento del fluido. Para simplificar, supondremos que $\Omega$ es independiente del tiempo.

\subsection*{1.1.1\quad Descripción del flujo}

Existen dos formas de describir el movimiento del fluido: la formulación \textit{lagrangiana} y la \textit{euleriana}. Utilizaremos aquí la descripción euleriana, que se basa en determinar la velocidad
\[
\bm{v}(x,t) = (v_1(x,t), \dots, v_N(x,t))
\]
de la partícula de fluido que pasa por el punto $x$ en el instante $t$.

Adoptaremos la siguiente notación:
\begin{itemize}
  \item $\rho$: densidad,
  \item $p$: presión,
  \item $\theta$: temperatura absoluta.
\end{itemize}
Estas cantidades se denominan \textit{variables de estado}.


\subsection*{1.1.2\quad Derivada material}

En la formulación euleriana, cada magnitud escalar $\phi = \phi(x,t)$ depende de la posición y del tiempo. Consideramos el cambio total de $\phi$ con respecto al tiempo a lo largo de la trayectoria de una partícula. Si $x = x(t)$ denota dicha trayectoria y $\phi = \phi(x(t),t)$, entonces:

\[
\frac{d\phi}{dt} = \frac{\partial \phi}{\partial t} + \nabla \phi \cdot \frac{dx}{dt}.
\]

Como $\frac{dx}{dt} = \bm{v}$, se sigue que:

\begin{equation}
\frac{d\phi}{dt} = \frac{\partial \phi}{\partial t} + \bm{v} \cdot \nabla \phi.
\tag{1.1}
\end{equation}

Este operador se llama \textit{derivada material} o \textit{derivada sustancial} y se denota comúnmente por $D\phi/Dt$.

\subsection*{1.1.3\quad Ecuaciones de conservación}

Las ecuaciones que gobiernan el movimiento del fluido provienen de leyes de conservación: masa, cantidad de movimiento y energía. Se basan en los siguientes principios físicos:

\begin{itemize}
  \item La masa no se crea ni se destruye (ley de conservación de la masa).
  \item La tasa de cambio de la cantidad de movimiento en un volumen de control es igual a la suma de las fuerzas externas (segunda ley de Newton).
  \item La tasa de cambio de la energía total es igual al trabajo de las fuerzas externas más el calor suministrado.
\end{itemize}

Estas leyes se formulan primero en forma integral sobre un volumen de control arbitrario, y luego se convierten en forma diferencial bajo hipótesis de suavidad suficientes.
